% The pronunciation line follows docs/lang/fr.md: silent letters are simply
% absent, a liaison is carried across on a hyphen onto the word it has moved
% to, an elision makes one word and drops the apostrophe, a tilde is a nasal
% vowel, é and è are the closed and open e, o and ò the closed and open o, ou
% eu ü keep their French values, ʃ ʒ ɲ are the sounds spelt ch j gn, and no
% letter is ever doubled.  \vb gives infinitive, first person present, past
% participle, and in brackets after the meaning what those three cannot say:
% the auxiliary where it is être, the nous form and the future where they
% are irregular -- never the class, which the infinitive already prints (the
% brackets here said -er and -re once); a noun carries its gender.
%
% The narration is in the past historic, which is none of those three forms,
% so each one is named where it stands -- the convention a French edition
% turns on, and the reason this text was written in that tense.
%
% What this fixture exists to protect is the apostrophe: every chunk below
% carries an elision, and joined with single spaces they must reproduce the
% source character for character, apostrophes included.
\chapopen{1}

\parstart{1.1}{C'était l'hiver, et le vieil homme}
\parnum{1.1}
\begin{frank}
\ch{}{C'était l'hiver,}{sétè livèr}{\vb{être}{ètr}{suis}{süi}{été}{été}{to be (nous \pw{sommes} \textit{sòm}; fut. \pw{serai} \textit{seré})}; \pw{c'} is \pw{ce}; \dw{hiver}{ivèr} winter, m., \pw{l'} for \pw{le} before a vowel}{it was winter,}
\ch{}{et le vieil homme}{é le vyèy òm}{\dw{vieux}{vyeu} old, f. \pw{vieille} \textit{vyèy}, and \pw{vieil} before a vowel; \dw{homme}{òm} man, m.}{and the old man}
\ch{}{n'avait plus de bois.}{navè plü de bwa}{\vb{avoir}{avwar}{ai}{é}{eu}{ü}{to have (fut. \pw{aurai} \textit{oré})}; \pw{ne … plus} no more, no longer, \pw{n'} for \pw{ne}; \dw{bois}{bwa} wood, m.}{had no more wood.}
\end{frank}

\parnum{1.2}
\begin{frank}
\ch{}{Aujourd'hui encore,}{oʒourdüi ãkòr}{\dw{aujourd'hui}{oʒourdüi} today; \dw{encore}{ãkòr} still, again}{even today,}
\ch{}{il regardait l'arbre}{il regardè larbr}{\vb{regarder}{regardé}{regarde}{regard}{regardé}{regardé}{to look at}; \dw{arbre}{arbr} tree, m.}{he would look at the tree}
\ch{}{devant sa porte.}{devã sa pòrt}{\dw{devant}{devã} in front of; \dw{porte}{pòrt} door, f.}{in front of his door.}
\end{frank}

\parend

\parstart{1.2}{Un enfant s'approcha et lui demanda}
\parnum{2.1}
\begin{frank}
\ch{}{Un enfant s'approcha}{ẽ-nãfã sapròʃa}{\dw{enfant}{ãfã} child, m., the \textit{n} of \pw{un} carried onto it; \vb{s'approcher}{sapròʃé}{m'approche}{mapròʃ}{approché}{apròʃé}{to come near (aux. être); here \pw{s'approcha} \textit{sapròʃa}, past historic}}{a child came near}
\ch{}{et lui demanda}{é lüi demãda}{\vb{demander}{demãdé}{demande}{demãd}{demandé}{demãdé}{to ask; here \pw{demanda} \textit{demãda}, past historic}}{and asked him}
\ch{}{ce qu'il attendait.}{se kil atãdè}{\pw{qu'} is \pw{que}; \vb{attendre}{atãdr}{attends}{atã}{attendu}{atãdü}{to wait for}}{what he was waiting for.}
\end{frank}

\parnum{2.2}
\begin{frank}
\ch{}{Le vieil homme sourit}{le vyèy òm souri}{\vb{sourire}{sourir}{souris}{souri}{souri}{souri}{to smile; here \pw{sourit} \textit{souri}, past historic, spelt as the present \pw{il sourit} and said like the participle}}{the old man smiled}
\ch{}{et répondit}{é répõdi}{\vb{répondre}{répõdr}{réponds}{répõ}{répondu}{répõdü}{to answer; here \pw{répondit} \textit{répõdi}, past historic}}{and answered}
\ch{}{qu'il attendait le printemps.}{kil atãdè le prẽtã}{\dw{printemps}{prẽtã} spring, m., the season}{that he was waiting for the spring.}
\end{frank}

\chapend
